%% evenend-sample-box.tex：evenendの見本（tcolorboxの中の段組）
%% マニュアル（evenend-ja.tex）が貼り込む
\documentclass[paper=b5,fontsize=9pt]{jlreq}
\usepackage[columns=1]{evenend}
\usepackage[most]{tcolorbox}
\setlength\columnseprule{0.4pt}
\newcommand\para{吾輩は猫である．名前はまだ無い．どこで生れたかとんと見当がつかぬ．何でも薄暗いじめじめした所でニャーニャー泣いていた事だけは記憶している．\par}
\begin{document}
\section{箱の中の段組}
\begin{tcolorbox}[title=2段（multicols）]
箱の中の前の段落．前後の段落とのベースライン間は行送りのまま．
\begin{multicols}{2}
\para\para
\end{multicols}
箱の中の後の段落．
\end{tcolorbox}
\begin{tcolorbox}[title={3段，段間1字，罫なし（evenendmulticols）}]
\begin{evenendmulticols}[columnsep=1\zw,columnseprule=0pt]{3}
\para\para\para
\end{evenendmulticols}
\end{tcolorbox}
\begin{tcolorbox}[title=columnbreakと脚注]
\begin{multicols}{3}
\para\columnbreak
脚注を付ける\footnote{箱の中の脚注は箱の下に置く．}．\para\columnbreak
\para
\end{multicols}
\end{tcolorbox}
\end{document}
